% ECONOMICS_PROBLEM: {"id": "D31-458", "title": "Reconciling income-inequality measurement", "jel_code": "D31", "source_id": "OP-0360"}
% !TeX program = xelatex
\documentclass[11pt,a4paper]{article}
\usepackage[margin=23mm]{geometry}
\usepackage{fontspec}
\setmainfont{Latin Modern Roman}
\usepackage{amsmath,amssymb,mathtools}
\usepackage{xcolor,enumitem,booktabs,longtable,array,xurl,hyperref}
\definecolor{muted}{HTML}{53636C}
\hypersetup{colorlinks=true,urlcolor=blue,linkcolor=blue}
\setlength{\parindent}{0pt}
\setlength{\parskip}{5pt}
\newcommand{\R}{\mathbb R}
\newcommand{\E}{\mathbb E}
\newcommand{\N}{\mathbb N}
\newcommand{\ind}{\mathbf 1}
\newcommand{\Var}{\operatorname{Var}}
\newcommand{\Cov}{\operatorname{Cov}}
\newcommand{\supp}{\operatorname{supp}}
\DeclareMathOperator*{\argmin}{arg\,min}
\DeclareMathOperator*{\argmax}{arg\,max}
\newcommand{\entrymeta}[3]{{\sffamily\small\color{muted}\textbf{\texttt{#1}}\quad|\quad #2\par #3\par}}
\newcommand{\Problemref}[1]{\texttt{#1}}
\newcommand{\JELref}[2]{\texttt{#2} (JEL \texttt{#1})}
\begin{document}
\section*{Reconciling income-inequality measurement}
\textbf{Problem ID:} \texttt{D31-458}\quad \textbf{JEL:} \texttt{D31}\par
% BEGIN ECONOMICS STATEMENT
\label{problem:OP-0360}
{\sffamily\small\color{muted}Original subject group: Inequality and economic mobility\par}
\label{definitions:problem:30007642}
\textbf{Audit classification:} Background; proposed extension. Publisher metadata and abstract checked. The source explains concept and allocation disagreements; no precise mathematical identification problem is declared open there.
\subsection*{Definitions and precise research target}
\textbf{Complete editorial problem.} Fix a country, accounting period, adult population, and a declared income concept \(\theta\). The concept specifies capital gains, retained business earnings, imputed government services, household sharing, taxes, and transfers. Let \(y_i(\theta)\) be person \(i\)'s income under those definitions, and let \(F_\theta\) be its population distribution. Observed sources comprise tax records \(D_T\), household surveys \(D_S\), and aggregate national accounts \(D_N\), each with stated coverage and error processes. For \(p\in(0,1)\), define the top-share functional \(Q_p(F)=\int_{1-p}^{1}F^{-1}(u)\,du/\int_0^1F^{-1}(u)\,du\), assuming a positive finite mean. The measurement task is to identify or bound \(Q_p(F_\theta)\) and reconcile discrepancies across sources without treating an accounting allocation convention as observed household income. Quantify sensitivity to nonfilers, evasion, survey nonresponse, asset valuation, and income-sharing assumptions. Separate disagreements about \(\theta\) from statistical uncertainty conditional on \(\theta\). Construct a transparent set of estimates for a predeclared family of concepts. There need not exist one uniquely correct comprehensive measure; the research problem is reproducible measurement and defensible reconciliation, with remaining normative choices made explicit.

\textbf{Definitions and admissible scope.} Use a population probability measure on adults and an integrable real-valued income random variable \(y(\theta)\) with \(E[y(\theta)]>0\). The generalized inverse is \(F^{-1}(u)=\inf\{y:F(y)\ge u\}\); the quantile-integral definition of \(Q_p\) allocates mass at ties fractionally and remains valid with negative incomes, although the result can exceed one. The income concept \(\theta\) fixes units, timing, capital gains, owner-retained earnings, household splitting, taxes, and in-kind allocations. For a data law \(P_D\), define \(\mathcal F_\theta(P_D)\) as distributions consistent with all source-coverage, linkage, and error assumptions. The sharp top-share set is \(\{Q_p(F):F\in\mathcal F_\theta(P_D)\}\). Accounting totals constrain distributions but do not specify how unassigned income belongs to individuals. Here \(\theta\), population weights, \(p\), and the measurement/coverage constraints defining \(\mathcal F_\theta(P_D)\) are mathematical inputs. All admissible distributions obey the same accounting identities and stated bounds on linkage errors, unreported amounts, and missing-source allocations. The target is \([\inf_{F\in\mathcal F_\theta}Q_p(F),\sup_{F\in\mathcal F_\theta}Q_p(F)]\), or the full image set if it is disconnected; emptiness signals incompatible assumptions and data. For concept sensitivity choose a finite predeclared \(\Theta\) and report each conditional set separately as well as their union. Allocating national-account residuals requires an explicit admissible allocation family and is not identified by accounting totals alone.
\subsection*{Variants and assumption changes}
\begin{enumerate}
\item Fixed-concept comparison (editorial): hold \(\theta\) fixed and vary bounds on evasion and survey nonresponse. Report how the sharp top-share interval changes without attributing concept differences to statistical error.
\item Concept sensitivity (source-informed editorial): specify a finite family \(\Theta\), such as realized taxable income versus an accrual concept with business-retention allocations. Report \(\{Q_p(F):\theta\in\Theta,F\in\mathcal F_\theta(P_D)\}\) and keep normative allocation choices explicit.
\end{enumerate}
\subsection*{References and source audit}
\textbf{Support and access.} Publisher metadata and abstract checked. The source explains concept and allocation disagreements; no precise mathematical identification problem is declared open there. \textbf{Locator:} Publisher abstract and citation; JEP 39(2), pp. 103--126.
\begin{enumerate}\raggedright
\item\hypertarget{definitions:source:30007642:1}{} Conor Clarke; Wojciech Kopczuk. 2025. \emph{Measuring Income and Income Inequality}. Publisher abstract and citation; Journal of Economic Perspectives 39(2), 103--126.
\url{https://www.aeaweb.org/articles?id=10.1257/jep.20241424}.
DOI: \url{https://doi.org/10.1257/jep.20241424}.
\end{enumerate}

\subsection*{Common conventions: Source mathematical conventions}

These conventions apply to each entry unless explicitly replaced.

\textbf{Models and identification.} A primitive model \(m\) specifies all probability laws, feasible choices, information, equilibrium and measurement rules used by its target. The admissible class \(\mathcal M\) is fixed independently of outcomes. If \(P_O(m)\) is its observable probability law and \(\tau(m)\) the target, its identified set at observable law \(P\) is
\[
 \mathcal I_\tau(P)=\{\tau(m):m\in\mathcal M,\ P_O(m)=P\}.
\]
Point identification means that this set is a singleton. An empty set signals incompatibility of data and assumptions. Coordinate bounds need not describe the joint set. Bounds are sharp only if every retained target is attainable or arbitrarily approachable by compatible models. Finite bounds require explicit restrictions on unobserved quantities; unrestricted confounding, hidden wealth, extrapolation or arbitrary equilibrium selection can yield uninformative sets.

\textbf{Causal interventions.} A potential outcome \(Y_i(p)\) is the outcome under a complete policy regime \(p\), including timing, financing, information and market spillovers. The target population and units are fixed before treatment. With interference, use the complete assignment vector or a justified exposure mapping; individual no-interference notation alone cannot identify an economy-wide intervention. A difference of mean potential outcomes does not require the cross-world joint distribution. Conditioning on post-treatment migration, survival, enrollment or employment changes the estimand and needs separate justification.

\textbf{Statistical statements.} An estimator, confidence set, forecast or test specifies a sampling law, model class and loss/coverage criterion. Uniform \((1-\alpha)\)-coverage means \(\inf_{m\in\mathcal M}\Pr_m\{\tau(m)\in C_n\}\geq1-\alpha\), or an explicitly asymptotic version. The whole parameter space trivially has coverage, so informativeness requires a stated width, risk or power objective. A forecasting improvement is relative to a named benchmark, information set and evaluation population.

\textbf{Equilibrium and optimization.} An equilibrium comprises feasible plans, optimality, beliefs where needed, budget constraints and all market/resource clearing. The primitives or a stated benchmark define each correspondence \(\mathcal E(p;m)\). Nonemptiness is assumed or proved, never inferred just from compact policy sets. A selected-equilibrium target uses a declared selection rule; a robust target quantifies over all permitted equilibria. Use suprema or infima unless attainment is established. An \(\varepsilon\)-optimal policy is feasible and within \(\varepsilon\) of the finite optimum value. Report an empty feasible set separately.

\textbf{Welfare and accounting.} Normative utility, interpersonal normalization, social weights, numeraire, discounting and the use of public receipts are primitives. Behavioral utility need not coincide with the welfare criterion. Household welfare includes ownership income and rents once; adding profits again to compensating variation generally double-counts them. A monetary equivalent is defined by an explicit transfer and price convention, with monotonicity and range conditions. Average effects, mechanism contributions, transfers and welfare losses are different targets. Infinite sums require convergence and dynamic feasible plans require terminal or no-Ponzi conditions.

\textbf{Source versions.} A working-paper date, revision date and journal publication year are distinguished. A DOI for a journal version is not automatically the DOI of an earlier manuscript. Locators refer to the stated version and printed pages where specified. A metadata-only check does not verify a claim about a full-text conclusion. Every entry's source subsection narrows the provenance claim accordingly.
% END ECONOMICS STATEMENT
\end{document}
